\scanpage{18}
the latter, if $\operatorname{card}J=\operatorname{card}I-1=n-1$, we shall see (by induction) that the cardinality of the set of classes is $(n-1)+1=n$.

This being so, we see first that $E_i$ is in one \inserted{(unique)} \unsure{of|of} these $n$ classes, which are open sets. Let $U_{J,i}$ be the class (or ``component'') which contains $E_i$, and let us consider the two components $\Sigma_i^+,\Sigma_i^-$ of $X\setminus E_i$.

\noindent\begin{tabular}{@{}|p{.91\linewidth}@{}}
\uline{Lemma} Every class of $\Pi_{J\deleted{\unsure{i|j|,}}}^{\unsure{\circ|0}}$, distinct from $U_{J,i}$, is contained in $\Sigma_i^+$ or $\Sigma_i^-$
\end{tabular}

This is almost a tautology: To say \deleted{the} \inserted{that a} given class \inserted{$U_{\unsure{\alpha|x}}$\unsure{,} of $\Pi_J^\circ$} does not contain $\deleted{\unsure{\Sigma}}\unsure{\overline{E_i}|E_i}$ i.e. is $\ne U_{J,i}$, means that there exists a $j\in J$ such that the $\Sigma_j^+\in\pi_0(X\setminus \deleted{\unsure{\Sigma}}E_j)$ \deleted{\unsure{chosen|chosen}} ``\unsure{component|component}'' of this class \deleted{is} is the element of the $\pi_{\unsure{0}}$ which does not contain $E_i$; it is therefore contained in one of the elements \inserted{(say $\Sigma_i^-$)} of $\pi_0(X\setminus\unsure{\Sigma_i|\overline{\Sigma_i}})$, by the study of the case \deleted{\unsure{two}} \inserted{$\operatorname{card}I=2$}, \deleted{\unsure{OK|this|it}} \unsure{a|a} fortiori $U_\alpha\subset\Sigma_i^-$, OK.
% archive sheet number: 17
